
	\newpage
% =========================================================	
	\section{Matrixnorm} \label{chap:matrixnorm}
% =========================================================

	Eine Matrix kann den relativen Eingabefehler beeinflussen. Um Fehler bei Vektoren und Matrizen messen zu können, brauchen wir Normen.
	
	\sbreak
	\subsection{Euklidische Norm} \label{chap:euklidischenorm}
		
		\defi{
			\n
			Die euklidische Norm wurde früher schon einmal beigebracht, deswegen nur eine kurze Wiederholung. Diese Norm gilt nur für Vektoren und entspricht der Länge eines Vektors:
			\begin{align*} 
		 		\norm{x}_2 = \sqrt{x_0^2+ \dots + x_n^2} 
		 	\end{align*}
			So wissen wir zum Beispiel, wenn wir mit einem Vektor der Länge eins normiert rechnen, eine Operation ausführen die die Länge nicht ändert und diese trotzdem nicht mehr 1 beträgt, dass sich ein Arithmetik Fehler eingeschlichen haben muss.
		}
	
	\sbreak
	\subsection{Spektralnorm} \label{chap:spektralnorm}
		
		\defi{
			\n
			Oder auch Matrix-Grenzen-Norm-2 genannt, beschreibt folgende Normierung für beliebige Matrizen:
			\begin{align} 
		 		\norm{A}_2 = \sup_{x \ne 0} \frac{\norm{Ax}_2}{\norm{x}_2} 
		 	\end{align}
		}
	
	\sbreak
	\methode{Kleiner Trick:
		\n
		Das werden wir zwar nicht beweisen, aber folgende Gleichungen gelten:
		\begin{align}
			\norm{A}_2 = \sqrt{\lambda_{\max}(A\trans A)}
		\end{align}
		Für symmetrische Matrizen:
		\begin{align}
			\norm{A}_2 = \max_{i\in[n]} |\lambda_i|
		\end{align}
		\n Für Diagonalmatrizen:
		\begin{align}
		\norm{A}_2 = \max_{i\in[n]} |A_{i,i}|
		\end{align}
		Eigenwert der Inversen:
		\begin{align}
			\lambda_i(A) \ne 0 \quad \forall i \Longrightarrow \lambda_i(A\inv) = \frac{1}{\lambda_i}
		\end{align}
	}
	
%		Für eine Vertiefung eignet sich Wikipedia: \\
%		}
	
	\sbreak
	\subsection{Rechen Eigenschaften}
	
	\defi{
		\n 
		Mit Normierungen müssen wir natürlich auch rechnen können, deswegen hier eine schnelle Übersicht über alle Eigenschaften.	
		\n
		Mit Skalar $a$, Vektor $v$, Matrix $A$,$B$ und $X$,$Y$ als Vektor oder Matrix gilt
		\begin{align*}
			\norm{aX}_2 &= \norm{a} \cdot \norm{X}_2 \\
			\norm{X+Y}_2 &\le \norm{X}_2 + \norm{Y}_2 \\
			\norm{A \cdot B}_2 &\le \norm{A}_2 \cdot \norm{B}_2 \\
			\norm{A \cdot v}_2 &\le \norm{A}_2 \cdot \norm{v}_2
		\end{align*}
	}


	\newpage
	\subsection{Kondition einer Matrix}
	
	\defi{
		\n 
		Mit der Einführung von Normen können wir nun auch noch die Kondition bzw. die Konditionszahl (siehe \refChapter{chap:kondition}) einer invertierbaren Matrix $A$ bestimmen.
	}

	\sbreak
	\methode{
		\n 
		Die Kondition der Matrix $A$ bezüglich der euklidischen Norm ist gegeben durch
		\begin{align}
			\cond(A) = \norm{A\inv}_2 \cdot \norm{A}_2
		\end{align}
		Außerdem kann man die Kondition der Matrix $A$ berechnen, indem man den größten Singulärwert der Matrix durch den kleinsten teilt.
		\n
		\fat{Falls die Matrix symmetrisch positiv definit} ist, kann man obiges vereinfachen zu:
		\begin{align}
			\cond(A) = \abs{\frac{\lambda\umax(A)}{\lambda\umin(A)}}
		\end{align}
		wobei $\lambda$ hierbei Eigenwerte beschreiben. $\lambda\umax(A)$ ist also der größte Eigenwert der Matrix $A$, $\lambda\umin(A)$ der kleinste.
	}


	\newpage
	\subsection{Kondition eines LGS}
	
	\defi{
		\n
		Ein LGS hat die Form $Ax = b$. Auch hiervon können wir die Kondition mithilfe der Normen bestimmen. Wir definieren $\tilde{x} := x + \delta x$
		als eine gestörte bzw. Näherungslösung des LGS. $\tilde{b} := b + \delta b$ 
		ist eine gestörte rechte Seite des LGS. %, während $\tilde{A}% := A + \delta A$ 
		%$ eine verfälschte Matrix beschreibt. %Die $\delta$ Variablen sind jeweils die Änderungen zu den exakten Werten.
	}

	\sbreak
	\methode{ Störung der rechten Seite
		\n 
		Sei die rechte Seite des LGS gestört, erhalten wir eine Näherungslösung:
		\begin{align}
			A\tilde{x} = \tilde{b}
		\end{align}
		Nun kann man die Fehler $e$ und $r$ beschreiben mit
		\begin{align}
			e &:= x - \tilde{x} \\
			r &:= b - \tilde{b} = Ae
		\end{align}
		Dann gilt:
		\begin{align}
			\frel &= \frac{\norm{e}}{\norm{x}} \\
			\frac{\norm{r}}{\norm{b} \cdot \cond(A)} & \le \frac{\norm{e}}{\norm{x}} \le \cond(A)\frac{\norm{r}}{\norm{b}}
		\end{align}
	}

%		Sei die Matrix $A$ verfälscht, gilt:
%	}